
\documentclass[11pt]{article}

\usepackage{amsmath}
\usepackage{amsfonts}

\newcommand{\universe}{\mathcal{U}}
\newcommand{\nats}{\mathbb{N}}

\begin{document}

\noindent
HoTT book: Exercise 1.11.  Prove for any type $A$ we
have $\neg\neg\neg A \to \neg A$.  In HoTT $\neg A$ is defined to be
$A \to 0$.  So the left-hand side is
$$
   ((A \to 0) \to 0) \to 0
$$
Assume this is true (inhabited),
then an inhabitant is a function $f : ((A \to 0) \to 0) \to 0$.
But the only function having codomain $0$ is $0 \to 0$.
So we infer $0 \equiv (A \to 0) \to 0$.
But that means $(A \to 0) \to 0$ is not inhabited, hence $A \to 0$ is inhabited.
And that is $\neg A$.

\medskip\noindent
HoTT book: Exercise 1.15.  Prove that identity of indiscernables (IoI) follows
from path induction.  IoI is: for every $C : A \to \universe$
there is a function
\begin{equation} \label{eq:f}
   f : \prod_{x, y : A} \prod_{p : (x =_A y)} C(x) \to C(y)
\end{equation}
satisfying
\begin{equation} \label{eq:f-too}
   f(x, x, \text{refl}_x) = \text{id}_{C(x)}
\end{equation}
Proof: we do have $x$ and $y$ varying freely in \eqref{eq:f} so it is
valid to use path induction, that is, we can only look at the case $y =_A x$
and $p = \text{refl}_x$ in which case \eqref{eq:f}
becomes $f(x, x, \text{refl}_x)$ is an inhabitant of $C(x) \to C(x)$,
and we know that \eqref{eq:f-too} is one such.  QED

\medskip\noindent
HoTT book: Exercise 1.16.  Prove that addition in $\nats$ is commutative,
where $\nats$ are defined in Section~1.9 of the HoTT book.  That section
defines $\text{add} : \nats \to \nats \to \nats$ satisfying the definitional
equalities
\begin{subequations}
\begin{align}
   \text{add}(0, n) & \equiv n
   \label{eq:add-base}
   \\
   \text{add}(\text{succ}(m), n) & \equiv \text{succ}(\text{add}(m, n))
   \label{eq:add-induction}
\end{align}
\end{subequations}
We prove by mathematical induction (the recursor for $\nats$).  In the
base case we have $\text{add}(0, 0) \equiv 0$ by plugging
into \eqref{eq:add-base}.  Then for the induction step we assume
\begin{equation} \label{eq:commutivity-induction}
   \text{add}(m, n) \equiv \text{add}(n, m)
\end{equation}
from which \eqref{eq:add-induction} gives
$$
   \text{add}(\text{succ}(m), n) \equiv \text{succ}(\text{add}(n, m))
$$
from which we conclude
\begin{equation} \label{eq:add-induction-induction}
   \prod_{m : \nats} \text{add}(m, n) \equiv \text{add}(n, m)
\end{equation}
So now we do another mathematical induction induction with base case
$$
   \prod_{m : \nats} \text{add}(m, 0) \equiv \text{add}(0, m)
$$
derived by plug-in and induction step where we assume
\eqref{eq:add-induction-induction} and derive
$$
   \prod_{m : \nats} \text{add}(m, \text{succ}(n)) \equiv
   \text{succ}(\text{add}(n, m))
$$
from which we conclude
$$
   \prod_{n : \nats}
   \prod_{m : \nats} \text{add}(m, n) \equiv \text{add}(n, m)
$$
QED

\end{document}
